%!TEX TS-program = xelatex
%!TEX encoding = UTF-8 Unicode
% Gilberto G. Briscoe-Martinez — Industry Resume (Deedy Format)
% Based on the Deedy Resume template (Debarghya Das, Apache 2.0)
% via LimHyungTae/Awesome-PhD-CV submodule (see cv-templates/).
%
% Compile from the repository root with:
%   make resume
% or, from this directory, with:
%   TEXINPUTS=.:../cv-templates/Awesome-PhD-CV/deedy-format: xelatex resume.tex
%
% Philosophy: one page, evidence only. The cover letter carries the
% narrative; the resume carries the numbers, venues, and systems.
%
% Bullet pattern: [bold plain-language catch] (venue) — what it does, in
% terms an industrial-robotics reader knows: numbers. No research jargon
% in the catch; 2–3 lines max; no terminal periods.
%
% Search for "TODO" before sending: those are places where a number or
% a full title would strengthen a claim and only you know the value.
%
% Layout note: the class's \namesection leaves \centering active at top
% level, so anything placed outside the two minipages will be centered
% with a stretched last line. Keep all body content inside the minipages.
% Dates use a literal en dash (–) because the class's \location and
% \setmainfont lack Mapping=tex-text, so "--" prints as two hyphens.

\documentclass[]{deedy-resume-openfont}
\usepackage{fancyhdr}
\usepackage{fontawesome5}
\usepackage{pifont}
\newcommand{\githubstars}[1]{[{\faGithubSquare}\,\ding{72}\,\textbf{#1}]}

% -----------------------------------------------------------------------------
% Font weight overrides
%
% The Deedy class ships with Lato-Lig (Light) as the main font and
% Raleway-ExtraLight for the sans family, which read as very thin on screen.
% We bump both to medium-weight variants so the resume looks sharper while
% keeping the original Deedy layout intact.
%
% Available Lato weights in the submodule:
%   Lato-Hai (Hairline) / Lato-Lig (Light) / Lato-Reg (Regular)
%   Lato-Bol (Bold)     / Lato-Bla (Black)
% Available Raleway weights: Thin, ExtraLight, Light, Regular, Medium,
%   SemiBold, Bold, ExtraBold, Heavy.
% -----------------------------------------------------------------------------
\setmainfont[
  Path = fonts/lato/,
  BoldItalicFont = Lato-BolIta,
  BoldFont       = Lato-Bol,
  ItalicFont     = Lato-RegIta,
]{Lato-Reg}
\setsansfont[
  Scale   = MatchLowercase,
  Mapping = tex-text,
  Path    = fonts/raleway/,
]{Raleway-Medium}

% Bump the section heading + name-section weights so they pop more.
% (Mirrors the original macros from deedy-resume-openfont.cls but with
% heavier Lato/Raleway variants.)
\usepackage{titlesec}

% Section headings (\section{...}) -- use Bold Lato instead of Light.
\titleformat{\section}{
  \color{headings}\scshape\fontspec[Path = fonts/lato/]{Lato-Bol}\fontsize{14pt}{17pt}\selectfont \raggedright\uppercase
}{}{0em}{}

% Name section -- bigger name than the original Deedy default (20pt -> 32pt),
% Light first name + Bold last name, contact line a touch larger,
% and a faint divider rule underneath.
\renewcommand{\namesection}[3]{
  \centering{
    \fontsize{32pt}{36pt}\selectfont
    \fontspec[Path = fonts/lato/]{Lato-Lig}\selectfont #1
    \fontspec[Path = fonts/lato/]{Lato-Bol}\selectfont #2
  } \\[4pt]
  \centering{
    \color{headings}\fontspec[Path = fonts/raleway/]{Raleway-Medium}\fontsize{10.5pt}{12pt}\selectfont #3}
  \noindent\makebox[\linewidth]{\color{headings}\rule{\paperwidth}{0.4pt}}
  \vspace{-12pt}
}

\pagestyle{fancy}
\fancyhf{}
% "Last updated" in the bottom-left footer (small Raleway-ExtraLight),
% replacing Deedy's default top-right \lastupdated textblock, which
% collided with the enlarged name.
\fancyfoot[L]{%
  \color{date}\fontspec[Path = fonts/raleway/]{Raleway-ExtraLight}\fontsize{8pt}{10pt}\selectfont
  Last updated \today%
}
\renewcommand{\footrulewidth}{0pt}

\begin{document}

%%%%%%%%%%%%%%%%%%%%%%%%%%%%%%%%%%%%%%
%     TITLE / CONTACT LINE
%%%%%%%%%%%%%%%%%%%%%%%%%%%%%%%%%%%%%%
\namesection{Gilberto G.}{Briscoe-Martinez}{
  \urlstyle{same}\href{https://gilberto.martinez.phd}{gilberto.martinez.phd} |
  \href{mailto:gilberto.briscoe-martinez@colorado.edu}{gilberto.briscoe-martinez@colorado.edu} |
  \href{https://www.linkedin.com/in/gbmartinez}{LinkedIn} |
  \href{https://github.com/ggb367}{GitHub} |
  \href{https://scholar.google.com/citations?user=6MNXSfsAAAAJ}{Scholar}
}

%%%%%%%%%%%%%%%%%%%%%%%%%%%%%%%%%%%%%%
%     COLUMN ONE (LEFT, 33%)
%%%%%%%%%%%%%%%%%%%%%%%%%%%%%%%%%%%%%%
\begin{minipage}[t]{0.33\textwidth}
\raggedright  % comma-separated skill lists wrap without justification stretch or hyphenation

%--------------- SKILLS ---------------
% ML first: this is what the reader is screening for.
\section{Skills}
\subsection{Machine Learning}
PyTorch, diffusion policies, imitation learning, ACT, VLA and VLM fine-tuning (LoRA), graph encoders (GAT, HGT), evaluation and ablation design, CUDA, NumPy
\sectionsep

\subsection{Robotics}
Motion planning, trajectory optimization, compliant and contact-rich control, non-prehensile manipulation, multi-robot coordination, ROS 2, Isaac Sim, OpenCV
\sectionsep

\subsection{Software}
Python, C++, CMake
% TODO: add Docker, AWS, or GCP here if you have used them; the JD lists them.
\sectionsep

\subsection{Platforms}
Franka Panda, UR5e, UR20, NVIDIA Jetson and Orin, industrial induction cells
\sectionsep

%--------------- EDUCATION ---------------
\section{Education}

\subsection{University of Colorado Boulder}
\descript{Ph.D.\ in Computer Science}
\location{Expected May 2027 | Boulder, CO}
HIRO Group \\
Advisor: Prof.\ A.\ Roncone
\sectionsep

\subsection{University of Texas at Austin}
\descript{B.S.\ in Aerospace Engineering}
\location{May 2021 | Austin, TX}
Human Centered Robotics Lab \\
(2018–2021)
\sectionsep

%--------------- AWARDS ---------------
\section{Awards}
\begin{tabular}{rl}
2022\,–\ 2026 & NASA NSTGRO Fellowship \\
           & (Grant 80NSSC22K1211) \\
\end{tabular}
\sectionsep

\end{minipage}
\hfill
%%%%%%%%%%%%%%%%%%%%%%%%%%%%%%%%%%%%%%
%     COLUMN TWO (RIGHT, 66%)
%%%%%%%%%%%%%%%%%%%%%%%%%%%%%%%%%%%%%%
\begin{minipage}[t]{0.66\textwidth}

%--------------- EXPERIENCE ---------------
\section{Experience}

\runsubsection{HIRO Group, CU Boulder}
\descript{| PhD Researcher \& NASA NSTGRO Fellow}
\location{Aug 2021 – Present | Boulder, CO}
\vspace{0.3cm}  % needed: the first list after a \section otherwise overlaps the location line
\begin{tightemize}
\item Built \textbf{DEFT}, a diffusion policy that recovers from any joint failure on a 7-DoF Franka arm without retraining: 99.5\% task success vs 42.4\% for RRT, 10/10 real-world tasks with multiple joints locked, trained on 4.7M trajectories (ICRA 2026)
\item Designed a fail-active controller for a two-arm package induction cell that reroutes work around faulty arms and conveyors instead of halting: a diffusion policy conditioned on a graph of each component's remaining capacity; 10$\times$ the throughput of the same policy without the graph under severe faults (ICRA 2027, under review)
% TODO: confirm the 10x against the final Table I numbers.
\item Tested whether the language model in a vision-language-action (VLA) policy affects the robot's actions, using logs from a production sorting arm; perturbing its output barely changed the picks, so the policy was mostly behavior cloning
\item Developed pushing and poking primitives that complete tasks after an end-effector failure: reachable workspace up to 79\% larger, 88.9\% task success (IROS 2024)
\item Co-developed \textbf{CERG}, a controller that bounds the arm's energy during contact (IFAC WC 2026, oral); mentor junior PhD students and undergraduates
\end{tightemize}
\sectionsep

\runsubsection{Plus One Robotics}
\descript{| Robotics AI/ML Researcher (Part-time)}
\location{June 2025 – Present | Boulder, CO}
\begin{tightemize}
\item Built the proof of concept for learned motion planning on the package induction cell, aimed at cutting new-site spin-up from 8–12 weeks to 6–10 by removing most post-install tuning; the CTO reassigned about 6 engineers to productionize it
\item Fine-tuned a VLA policy on 20 hours of test-cell data; after 3 months it runs the cell at 65\% of the classical system's throughput and retries failed picks without the hand-coded confirmation step the classical system needs
\item Designed the policy architecture the team builds on: a scalable multi-model task and motion planning approach, physically-grounded through a graph representation of the cell, so one policy runs across multiple cells with different arms and conveyors
\item Built a half-size test cell with robots and conveyors in a closed loop, and a five-conveyor, two-arm Isaac Sim model that has generated 150+ hours of training data; test on the full-scale cell in San Antonio every two months
\end{tightemize}
\sectionsep

\runsubsection{Human Centered Robotics Lab, UT Austin}
\descript{| Undergraduate Researcher}
\location{Jan 2018 – May 2021 | UT Austin Villa RoboCup@Home | Austin, TX}
\begin{tightemize}
\item Trained object-recognition models for the team's home-service robot (2019) and implemented data augmentation techniques that improved recognition performance from 60\% to >95\% in real world; co-authored the team's RoboCup 2019 report
% TODO: a dataset size or accuracy delta here would match the DEFT bullet's weight.
\item Built the robot's navigation, task planning, and behavior-tree task execution with Prof.\ Luis Sentis and Prof. Justin Hart; co-authored the team's AAAI AI-HRI 2019 report
\end{tightemize}
\sectionsep

\runsubsection{METECS, NASA Johnson Space Center}
\descript{| Robotics Software Intern}
\location{May 2020 – Aug 2020 | Dexterous Robotics Lab, Houston, TX}
\begin{tightemize}
\item Wrote affordance templates, reusable task models of grasp points and motions, so Robonaut 2, NASA's International Space Station humanoid, could open latches and doors
\end{tightemize}
\sectionsep

\runsubsection{Johns Hopkins APL}
\descript{| Robotics Software Engineer (Intern)}
\location{June 2021 – July 2021 | Remote, Laurel, MD}
\begin{tightemize}
\item Built a ROS 1 and ROS 2 simulator with Gazebo models of real buildings so the team could keep developing between tests on the physical robots
\end{tightemize}
\sectionsep

%--------------- PUBLICATIONS ---------------
% Replaces the old "Selected Projects" block, which duplicated the
% experience bullets. A real list reads as "published doctorate" at a
% glance, which is the HFC screen.
% TODO: swap short-form titles for full titles where marked, and add a
% marker (e.g. \dag) on first-author entries.
\section{Publications}
\vspace{0.3cm}
\begin{tightemize}
\item \textbf{ICRA 2027} (under review) — \textit{What Can My Neighbor Still Do? Fail-Active Control for Process-Coupled Robot Cells via Graph-Conditioned Diffusion}
\item \textbf{ICRA 2027} (under review) — \href{https://arxiv.org/abs/2603.14634}{PBD-R}: Rigid-Body Dynamics in Particle-Based Simulation % TODO full title if it fits
\item \textbf{ICRA 2026} — DEFT: diffusion policy for joint-failure recovery % TODO full title
\item \textbf{IFAC WC 2026} (oral) — CERG: Compliant Explicit Reference Governor % TODO full title
\item \textbf{IROS 2024} — task recovery after end-effector failure via pushing and poking % TODO full title
\end{tightemize}
Issued U.S.\ Patent \href{https://patents.google.com/patent/US10339827B2}{\textbf{10,339,827}}. Full list on \href{https://scholar.google.com/citations?user=6MNXSfsAAAAJ}{Google Scholar}.
\sectionsep

\end{minipage}
\end{document}